% ═══════════════════════════════════════════════════════════════════
% ЧАСТЬ IV. ТРЕТИЙ СТЕНД: ЯКОБИАН КВАРТИКИ КЛЕЙНА
% ═══════════════════════════════════════════════════════════════════
\part{Третий стенд: якобиан квартики Клейна}

\section{Квартика Клейна и её инварианты}\label{sec:klein}

Третья ступень лестницы — якобиан квартики Клейна
\[
  X\;:\; x^3y+y^3z+z^3x=0 \qquad \subset\ \mathbb{P}^2.
\]
Кривая имеет род $g=3$, группу автоморфизмов $\Gamma(2,3,7)$ порядка
$336$ (наибольшую среди кривых рода $3$) и три выделенных поворота
с фазами $\pi/2$, $\pi/3$, $\pi/7$. Якобиан $\Jac(X)$ изогенен кубу
эллиптической кривой с комплексным умножением
$\QQ(\sqrt{-7})$ — теорема Рибета; программа подтверждает все
целочисленные следствия этой структуры точной арифметикой.

\subsection{Гладкость и редукция к форме Вейерштрасса}

\begin{proposition}[гладкость и род]\label{prop:klein-smooth}
Квартика Клейна гладка: если $F_x=F_y=F_z=0$ на кривой, то
$28(xyz)^3=0$, значит $xyz=0$, и в каждом граничном случае система
силится в $(0,0,0)$ — не проективной точке. Род равен
$\tfrac{(4-1)(4-2)}{2}=3$.
\end{proposition}

\begin{proof}
Частные производные $F_x=3x^2y+z^3$, $F_y=x^3+3y^2z$, $F_z=y^3+3z^2x$.
Комбинация $xF_x+yF_y+zF_z=4(x^3y+y^3z+z^3x)=0$ на кривой не даёт
нового условия; исключение переменных из трёх уравнений производных
вместе с уравнением кривой даёт тождество $28(xyz)^3=0$
(проверено символьной арифметикой — вычет по модулю кривой).
Три случая $x=0$, $y=0$, $z=0$ разбираются подстановкой: каждый
силится в тройном нуле. Род вычисляется по формуле Плюккера для
гладкой плоской кривой степени $4$.
\end{proof}

\begin{proposition}[форма $49a1$ и кратные $\{-7, -15\}$]\label{prop:klein-49a1}
Квартика Клейна над $\ZZ$ имеет модель $49a1$ с дискриминантом
$\Delta=-343=-7^3$, коэффициентом $c_4=105$ и $j$-инвариантом
\[
  j \;=\; -3375 \;=\; -15^3 .
\]
Кратная форма: $W^2=v^3-35v-98$, корни $e_1=7$,
$e_{2,3}=\tfrac{-7\pm\sqrt{-7}}{2}$; разложение
$v^3-35v-98=(v-7)(v^2+7v+14)$ точное.
\end{proposition}

\begin{proof}
Формулы Сильвермана для квартик над $\ZZ$ вычисляются в целых числах:
$\Delta=-343$, $c_4=105$, и $j=c_4^3/\Delta=-15^3=-3375$ — все
значения точные. Переход к кратной форме: по модулю уравнения кривой
справедливо $(2y+x)^2=4x^3-3x^2-8x-4$ (символьное тождество);
замена $v=4x-1$, $W=4(2y+x)$ даёт $W^2=v^3-35v-98$. Разложение
кубического многочлена проверяется раскрытием скобок: $(v-7)(v^2+
7v+14)=v^3-35v-98$.
\end{proof}

\subsection{Период: три независимые схемы}

\begin{proposition}[вещественный период]\label{prop:klein-period}
Вещественный период $\Omega$ модели $49a1$, вычисляемый как
$2\int_7^\infty dx/y$, согласован по трём независимым схемам:
прямая квадратура, замена $t^2$ и редукция к симметричному интегралу
Карлсона $\mathrm{RF}$:
\[
  \Omega \;=\; 1.93331170561681154673308\ldots
  \qquad (\text{согласие схем}\ <1.4\cdot10^{-32}\ \text{при}\ 60\ \text{цифрах}).
\]
\end{proposition}

\begin{proof}
Схема~1: прямое квадратурирование по $(7,\infty)$ с разрезом в
особенностях. Схема~2: замена $u^2=v-7$ снимает корень; интеграл
сводится к гладкой функции на $[0,\infty)$. Схема~3: редукция к
симметричному интегралу Карлсона $\mathrm{RF}(x,y,z)$ с точными
преобразованиями аргументов. Все три схемы не зависят друг от друга
ни алгоритмически, ни по коду; совпадение с точностью $10^{-32}$
при $60$ рабочих цифрах исключает программные совпадения.
\end{proof}

\subsection{Периодное отношение и «венец» стенда}

\begin{theorem}[восстановление $j$ из периода]\label{th:klein-j}
Из численного периода восстанавливается модульное отношение
$\tau=(1+\sqrt{-7})/2$ и $j$-инвариант $j(\tau)=-3375$ с точностью
$57$ знаков (якорь $j(i)=1728$). Тождество характера и периода
$\zeta+\zeta^2+\zeta^4=\tau-1$ точное; изотипика $H^0(\Omega^1)$
равна $(1,1,0,1)$ по характерам $\mu_4$; $\lambda_1$ Хивуда
$=\sqrt2$ кратности $6$.
\end{theorem}

\begin{proof}
Периодное отношение $\tau$ вычисляется по $\lambda$-модулю решётки:
ангармоническая группа $S_3$ переставляет шесть значений $\lambda$;
$K$-формула с $\mathrm{SL}_2$-нормировкой выделяет каноническое;
отклонения $<10^{-30}$. Подстановка $\tau=(1+\sqrt{-7})/2$ в
$q$-разложение $j$ с якорем $j(i)=1728$ (глубина ряда до
$q^{-1}\sim e^{-2\pi}$) даёт $57$ согласованных знаков значения
$-3375$. Тождество характера: сумма Гаусса $\zeta_7+\zeta_7^2+
\zeta_7^4$ равна $(-1+\sqrt{-7})/2=\tau-1$ — проверяется возведением
в квадрат обеих частей в кольце $\ZZ[\zeta_7]$. Изотипика: действие
$\mu_4$ на трёхмерном пространстве голоморфных форм раскладывается
$(1,1,0,1)$ — проекционные операторы вычислены точно.
$\lambda_1=\sqrt2$ кратности $6$: собственные значения оператора
Хивуда считаны точным характеристическим многочленом.
\end{proof}

\begin{databox}[Клейн, сводка стенда]
Гладкость: $28(xyz)^3=0$ (точное тождество); род $3$;
$\Delta=-343=-7^3$; $c_4=105$; $j=-3375=-15^3$; кратная форма
$W^2=v^3-35v-98$; корни $7,\ \tfrac{-7\pm\sqrt{-7}}2$;
$\Omega$ по трём схемам с согласием $1.4\cdot10^{-32}$;
$\tau=(1+\sqrt{-7})/2$; $j(\tau)=-3375$ из периода ($57$ знаков);
$\zeta+\zeta^2+\zeta^4=\tau-1$; изотипика $(1,1,0,1)$;
$\lambda_1(\text{Хивуд})=\sqrt2$, кратность $6$;
код Фано: $C_7+\mu_4$.
\end{databox}

Стенд Клейна замыкает тройку калибровочных стендов: на нём фаза
$\pi/7$ получает арифметическое окружение — гептаду $\sqrt{-7}$,
поле комплексного умножения и точный $j$-инвариант. Все три величины
$(n,B,\lambda_0)=(7,22,\lambda_1)$ согласованы сертификатами~B и~C
на этой ступени, что и позволяет переносить конвейер на верхние
уровни лестницы.
